\documentclass[planificacion2024.tex]{subfiles}
\usepackage{comment,setspace}

\setstretch{1.5}
\begin{document}

\section*{DIAGNÓSTICO AULICO: AÑO 2024}
\begin{tabular}{r l}
\textbf{CURSO :}& 4° 1°\\
\end{tabular}

\vspace{1cm}
Hay varios aspectos importantes sobre su desempeño académico y su comportamiento en clase. En primer lugar, es evidente que estos estudiantes poseen una sólida base en matemáticas, lo cual es fundamental para su formación en Electrotecnia II. Su capacidad para comprender y aplicar conceptos matemáticos les proporciona una ventaja significativa en el ámbito académico y en futuras actividades profesionales relacionadas con su especialidad.

Además, estos alumnos muestran una actitud positiva hacia el trabajo y una buena predisposición para aprender. Durante las clases, se observa que prestan atención y participan activamente en las actividades propuestas. Esta disposición favorable contribuye a crear un ambiente de aprendizaje dinámico y colaborativo, donde se fomenta el intercambio de ideas y el trabajo en equipo.

Sin embargo, a pesar de estas fortalezas, los estudiantes enfrentan varios desafíos que afectan su rendimiento académico. Uno de los problemas más destacados es la presencia de distracciones en el aula. El uso excesivo del celular durante las clases y el hábito de llevar objetos personales como cartas son prácticas que interfieren con su concentración y comprometen su capacidad para seguir el ritmo de la enseñanza. Estas distracciones también se reflejan en la dificultad para completar las tareas asignadas en tiempo y forma, lo que impacta negativamente en su rendimiento académico y en su desarrollo personal.

Además, la falta de acceso a internet o su mal funcionamiento constituyen otro obstáculo importante. En un mundo cada vez más digitalizado, el uso de recursos en línea es fundamental para complementar el aprendizaje en el aula y proporcionar a los estudiantes una experiencia educativa enriquecedora. La falta de acceso a estos recursos limita sus oportunidades de aprendizaje y dificulta su capacidad para adquirir habilidades relevantes para su futura carrera.

A pesar de estos desafíos, es importante reconocer el potencial de estos alumnos y trabajar en conjunto para superar las dificultades que enfrentan. Mediante estrategias educativas innovadoras, apoyo individualizado y una estrecha colaboración de los estudiantes , es posible crear un ambiente de aprendizaje inclusivo y estimulante que promueva el éxito académico y el desarrollo integral de los alumnos.

\begin{comment}
prompt de chatGPT

Elaborar un diagnóstico áulico un texto de 200 palabras, en lenguaje de profesor de escuela, que muestre las capacidades de alumnos de un colegio técnico con orientación en electricidad que tienen buena base en matemáticas, y que tambien tienen buena predisposición ypara trabaja y que prestan atención, pero algunos alumnos tienen problemas distractivos de la atención como el uso del celular, llevan cartas y ponerse a trabajar luego no entregando los trabajos, cosa que debería estar prohibida en el colegio, pero dentro de todos son buenos alumnos. Además tambien es un problema la falta de internet o el mal funcionamiento de la misma como recurso didáctico.


\end{comment}

\end{document}

